% Created 2023-09-10 Sun 20:34
% Intended LaTeX compiler: pdflatex
\documentclass{moderncv}
\usepackage[utf8]{inputenc}
\usepackage[T1]{fontenc}
\usepackage{graphicx}
\usepackage{longtable}
\usepackage{wrapfig}
\usepackage{rotating}
\usepackage[normalem]{ulem}
\usepackage{amsmath}
\usepackage{amssymb}
\usepackage{capt-of}
\usepackage{hyperref}
\usepackage[scale=0.75]{geometry}
\usepackage[T1]{fontenc}
\usepackage{lmodern}
\usepackage[sorting=ydnt,minnames=5,minbibnames=5,maxbibnames=5,backref=true]{biblatex}
\usepackage{multicol}
\moderncvstyle{classic}
\moderncvcolor{blue}
\name{Fred K. Gruber}{}
\email{fgruber@ieee.org}
\phone[mobile]{(+1) 617-784-1265}
\homepage{www.fredgruber.org}
\address{Newton, MA}
\social[linkedin]{fredgruber}
\date{\today}
\title{Senior Principal Scientist}
\hypersetup{
 pdfauthor={Fred K. Gruber},
 pdftitle={Senior Principal Scientist},
 pdfkeywords={},
 pdfsubject={},
 pdfcreator={Emacs 29.1 (Org mode 9.6.6)}, 
 pdflang={English}}
\begin{document}

\maketitle


\section{Summary}
\label{sec:org58f357a}
With 10 years of experience in analyzing medical datasets from clinical and real-world sources, I have a proven track record in the investigation of prognostic and predictive biomarkers as well as drug target discovery. I am proficient in all stages of data analysis, from preprocessing to modeling and interpretation, and have successfully led multiple projects with top pharmaceutical scientists as stakeholders. My expertise lies in statistical and machine learning models applied to medical data, with a focus on causal methods, high-dimensional variable selection, data integration, and interpretability. I am skilled in working with multi-modal high-dimensional genomic data and clinical variables, utilizing techniques such as bayesian networks, logistic regression, survival analysis, neural networks, and more. Additionally, I am the developer and maintainer of multiple internal R packages for data preprocessing, visualization, machine learning, and analysis as well as maintainer of the simcausal R package for simulation of causal networks.

\section{Professional Experience}
\label{sec:orgc2abaf0}
\cventry{\textbf{Aug 2018 -- current}}{Senior Principal Scientist}{Aitia}{Somerville, MA}{}{\begin{itemize}
\item Analysis of high-dimensional multi-modal biological data.
\item Experience with RNAseq, simple nucleotides variants, structural variants, copy number variants, and clinically relevant variables.
\item Focus on discovering biomarkers and drug targets for different diseases.
\item Experience in analyzing oncological datasets for lung cancer, multiple myeloma, prostate cancer, pancreatic cancer, and kidney cancer.
\item Utilization of different data sources such as TCGA/GDC, Flatiron, ICGC, Tempus MMRF, and others.
\item Involvement in all stages of the analysis process, from data preprocessing to model construction and interpretation as well as presentation to clients.
\item Use of proprietary causal Bayesian network learning and predictive tools as main tools.
\item Exploration of open source machine learning techniques applicable to datasets.
\item Investigation of unsupervised machine learning methods for tumor subclonal composition analysis and patient drug response stratification.
\item Developed the best performing model predicting high-risk multiple myeloma in the Multiple Myeloma DREAM Challenge finding a novel expression-based risk marker, PHF19, which was further validated in vitro and published in journal Leukemia.
\end{itemize}
}
\cventry{\textbf{Oct 2016 -- Oct 2018}}{Principal Scientist}{Aitia}{Cambridge, MA}{}{\begin{itemize}
\item Responsible for the development of causal network models for multiple myeloma based on multiple iterations of the CoMMPass MMRF  dataset.
\item Developed preprocessing pipelines to prepare diverse data modalities including copy number variants, structural variants, single nucleotide variants, gene expression, and clinical data.
\item Validation of result by literature search.
\item Published multiple peer-reviewed conference papers in a major blood cancer conference.
\end{itemize}
}
\cventry{\textbf{Jul 2013 -- Oct 2016}}{Senior Systems Biologist}{Aitia}{Cambridge, MA}{}{\begin{itemize}
\item Developed, implemented, and benchmarked algorithms for the prediction of rare adverse events on patients taking the client's drug.
\item Developed novel algorithms implementing variations of logistic regression that perform better under imbalanced outcomes.
\item Comparative analysis using 17 public-domain datasets
\item Applied to the client's clinical and genomic multiple sclerosis dataset.
\item Successfully delivered R package with new algorithm.
\item Success of algorithm lead to production code in the internal causal and predictive tools at Aitia.
\item Developed predictive and causal models of Chron's disease for the purpose of elucidading the relationships between TL1A, DR3,	IFNgamma, and a	variety	of genetic, gene expression, immune response, and clinical variables from tissue samples.
\end{itemize}
}

\cventry{\textbf{May 2009 -- Jun 2013}}{Postdoctoral Scientist}{Schlumberger-Doll Reasearch Center}{Cambridge, MA}{}{\begin{itemize}
\item Designed, implemented, and tested algorithms for analyzing Nuclear Magnetic Resonance measurements for oil exploration
\item Utilized Matlab extensively for research purposes
\item Utilized Python for production code integration with company software
\item Worked with optimization algorithms, sparse coding, regularization, and general machine learning tools.
\end{itemize}
}

\section{Education}
\label{sec:orgf4cad3d}
\cventry{\textbf{2004 -- 2009}}{Ph.D., Electrical engineering}{Northeastern University}{Boston, MA}{}{\begin{itemize}
\item Research Assistant at Northeastern University's Communications and Digital Signal Processing Center.
\item Member of the Honor Society of Eta Kappa Nu.
\item GPA of 3.9/4.0
\item Developed algorithms for the determination of the location of objects from the measure of scattered fields. This is a fundamental problem in the area of remote sensing and inverse problems. Also studied the fundamental limits of communication using electromagnetic fields.
\end{itemize}
}

\cventry{\textbf{2003 -- 2004}}{Master of science, Industrial Engineering and Management Systems}{University of Central Florida}{Orlando, FL}{}{\begin{itemize}
\item Research Assistant the Center for NASA Simulation Research.
\item Phi Kappa Phi honor society
\item GPA of 3.9/4.0
\item Recipient of the University of Central Florida's Graduate Merit Fellowship Award, 2004
\item Worked at the NASA Simulation Research Group and was involved in the design of an advanced simulation of the probability of failure a space shuttle during reentry and launch.
\item For my master thesis, I developed a Genetic Algorithm in C++ for the optimal selection of the parameters of a Support Vector Machine. Published this work in book chapter.
\end{itemize}
}

\cventry{\textbf{1998 -- 2003}}{Bachelors, Electrical and Electronic Engineering}{Universidad Tecnológica de Panamá}{Panamá City, Panamá}{}{\begin{itemize}
\item Member of IEEE
\item GPA of 3.74/4.0
\item Honor Scholarship for undergraduate studies in Engineering at the Technological University of Panama, 1998 - 2002.
\item Gold Medal "Guillermo Andreve" for educational achievements, Panama, 1997.
\end{itemize}
}

\section{Relevant Certificates}
\label{sec:org5952a1c}
\cventry{\textbf{2023}}{Combining information for Causal Inference}{Harvard T.H. Chan School of Public Health}{Cambridge, MA}{}{}

\cventry{\textbf{2022}}{Target Trial Emulation}{Harvard T.H. Chan School of Public Health}{Cambridge, MA}{}{}


\cventry{\textbf{2020}}{Applied Deep Learning Boot Camp}{MIT Professional School}{Cambridge, MA}{}{}

\cventry{\textbf{2018}}{An Introduction to Causal Inference}{Harvard T.H. Chan School of Public Health}{Cambridge, MA}{}{}

\cventry{\textbf{2016}}{Data Science: Data to insights}{MIT Professional School}{Cambridge, MA}{}{}

\section{Skills}
\label{sec:org3e62b7d}
\begin{multicols}{2}
\begin{itemize}
\item R
\item ggplot2
\item shiny
\item python
\item matlab
\item elisp
\item machine learning
\item causal discovery
\item causal inference
\item python
\item High-dimensional data analysis
\item AWS
\item data wrangling
\item data analysis
\item data visualization
\item statistics
\item teamwork
\item scientific presentation
\item research
\end{itemize}
\end{multicols}


\section{Languages}
\label{sec:org04911b8}
\begin{itemize}
\item \textbf{English} Fluent
\item \textbf{Spanish} Native
\end{itemize}


\begin{refsection}[resume.bib]
\nocite{*}
\bibliographystyle{ieee}
\printbibliography[title={Selected Publications}]
\end{refsection}
\end{document}